% =============================================================
% Fig 1: CoT-Faith workflow overview (TikZ skeleton — POLISH BY HAND)
% =============================================================
% Suggested placement: after \begin{document}, floated at [t] on p.1.
% Compile with pdflatex (loads tikz). Adjust node positions/colours
% to match the paper's colour scheme.
%
\begin{figure}[t]
\centering
\begin{tikzpicture}[
    node distance=0.7cm,
    box/.style={rectangle, rounded corners=2pt, draw=black, thick,
                minimum width=2.4cm, minimum height=0.9cm,
                align=center, font=\small},
    inbox/.style={box, fill=blue!8},
    modbox/.style={box, fill=green!12},
    probe/.style={box, fill=orange!12, minimum width=3.0cm},
    findbox/.style={box, fill=red!7, minimum width=3.4cm, minimum height=0.8cm},
    arr/.style={-{Latex[length=2mm]}, thick, shorten <=1pt, shorten >=1pt},
]

% Row 1: inputs
\node[inbox]  (img) at (0, 0)  {LIBERO image\\ ({\scriptsize$H\times W\times3$})};
\node[inbox, right=0.5cm of img]  (instr) {instruction\\ ({\scriptsize"pick up X"})};
\node[inbox, right=0.5cm of instr] (cot)   {\emph{GT} CoT\\ ({\scriptsize9-tag trace})};

% Row 2: model
\node[modbox, below=0.9cm of instr, minimum width=6.4cm, minimum height=1.0cm]
      (model) {CoT-VLA (ECoT-family or ablation variant)};

% Row 3: three probes
\node[probe, below=1.1cm of model.south west, xshift=0.6cm] (p1)
      {\textbf{Probe 1}: segment\\ attention (4 buckets)};
\node[probe, below=1.1cm of model] (p2)
      {\textbf{Probe 2}: causal\\ CoT edit ($\Delta_\infty$)};
\node[probe, below=1.1cm of model.south east, xshift=-0.6cm] (p3)
      {\textbf{Probe 3}: attention\\ $\to$ error AUROC};

% Row 4: findings
\node[findbox, below=0.9cm of p1] (f1)
      {attn shape = arch-\\determined (Fig 2)};
\node[findbox, below=0.9cm of p2] (f2)
      {reasoning target\\ = causal (Fig 3-4)};
\node[findbox, below=0.9cm of p3] (f3)
      {attn $\ne$ error\\ AUROC $\approx 0.5$};

% Arrows: inputs -> model
\draw[arr] (img)   -- ($(model.north)+(-1.8,0)$);
\draw[arr] (instr) -- (model.north);
\draw[arr] (cot)   -- ($(model.north)+(1.8,0)$);

% Model -> probes
\draw[arr] (model.south) -- (p1.north);
\draw[arr] (model.south) -- (p2.north);
\draw[arr] (model.south) -- (p3.north);

% Probes -> findings
\draw[arr] (p1) -- (f1);
\draw[arr] (p2) -- (f2);
\draw[arr] (p3) -- (f3);

\end{tikzpicture}
\caption{\textbf{CoT-Faith methodology}. Each LIBERO first-step
sample $(\mathrm{img}, \mathrm{instr}, \mathrm{cot}_\text{GT})$ is fed
to a CoT-VLA with the standard 9-tag ECoT prompt. Three probes measure
different aspects of chain-of-thought faithfulness: segment-decomposed
action-token attention (Probe 1); Δaction under causal edits along ten
semantic axes with a null selfsplice control (Probe 2); AUROC of
attention as a per-sample action-error predictor (Probe 3). Across
11 models we find that attention on the CoT segment is
architecture-determined, while causal effect requires reasoning-target
training — attention does not proxy causation.}
\label{fig:workflow}
\end{figure}
